\chapter{INTRODUCTION}%
\label{chap:introduction}

Embedded systems play a crucial role in modern electronic devices by integrating hardware and 
software to perform dedicated functions efficiently. Microcontroller-based systems are widely 
used in applications such as digital clocks, timers, control systems, and monitoring devices 
due to their low cost, reliability, and real-time processing capabilities. This project focuses 
on the design and implementation of a digital clock using an AVR microcontroller and 
seven-segment displays to demonstrate fundamental embedded system concepts.

\section{BACKGROUND AND MOTIVATION}
Digital clocks are one of the most common real-time embedded applications and serve as an 
excellent platform for learning microcontroller programming. The motivation behind this project 
is to gain practical exposure to embedded system design by interfacing a microcontroller with 
external display hardware and implementing precise time measurement using internal timers.

The AVR ATmega32 microcontroller provides built-in hardware timers, interrupt handling 
mechanisms, and configurable I/O ports, making it suitable for implementing timing-based 
applications. By developing a digital clock that displays time in minutes and seconds format, 
this project aims to strengthen understanding of timer configuration, interrupt-driven 
programming, and display multiplexing techniques.

\section{PROBLEM STATEMENT}
The objective of this project is to design and implement a digital clock capable of accurately 
displaying time in \textbf{MM:SS format}\nomenclature{MM:SS}{Minutes:Seconds} using an 
AVR ATmega32 microcontroller. The system must 
generate precise one-second intervals, update time values reliably, and display the output 
using four seven-segment displays while efficiently utilizing the available microcontroller I/O 
pins.

\section{PROJECT OBJECTIVES}
The main objectives of the project are:
\begin{itemize}
    \item To design a digital clock using the AVR ATmega32 microcontroller.
    \item To generate accurate one-second timing using hardware timers.
    \item To implement interrupt-driven time updating for reliable operation.
    \item To interface and control multiple seven-segment displays using multiplexing.
    \item To simulate and verify the system using Proteus software.
    \item To develop and test the embedded C program using Microchip Studio.
\end{itemize}
\enumitemEndSpacing{}

\section{SCOPE OF THE PROJECT}
The scope of this project is limited to the design and simulation of a digital clock displaying 
minutes and seconds using four seven-segment displays. The system uses the internal RC 
oscillator of the microcontroller for timing purposes and does not include external time 
calibration mechanisms.

The project does not cover advanced features such as hour display, alarm functionality, 
real-time clock (RTC)\nomenclature{RTC}{Real-Time Clock} modules, or power optimization 
techniques. However, the design can be 
extended in the future to incorporate additional features and hardware components.

\section{SIGNIFICANCE OF THE PROJECT}
This project is significant as it provides a practical understanding of real-time embedded 
system design using AVR microcontrollers. It demonstrates the effective use of timers, 
interrupts, and I/O port manipulation in a real-world application.

Additionally, the project serves as a foundational learning exercise for students studying 
embedded systems, as the concepts applied in this work are directly transferable to more 
complex applications such as industrial timers, digital control systems, and consumer 
electronic devices.

